\section{A renormalized endpoint trace and a controlled normalization}
\label{sec:spectral-renormalized-boundary}

The ordinary endpoint value in AC-S36 is zero.  A different,
precisely defined endpoint functional is available.  Its construction
uses the half-line Green formula for the geometric stable process;
it does not follow from a boundary upper bound alone.

\paragraph{Primary input and provenance.}
Grzywny--Ryznar, \emph{Potential theory of one-dimensional geometric
stable processes}, arXiv:1107.0745v1, \S3, pp.~5--7,
Proposition~3.2, Remark~3.6, and Lemma~4.3, p.~10,
give the renewal density and the half-line Green formula used below.
The formula in Lemma~4.3 is explicitly deduced there from Bertoin,
\emph{Levy Processes}, Theorem~20, p.~176; it is retained as a classical
fluctuation-theory input, not reproved here.
\url{https://arxiv.org/pdf/1107.0745}.
The renewal normalization is also compatible with Kwa\'snicki,
\emph{Spectral analysis of subordinate Brownian motions in half-line},
arXiv:1006.0524, Definition~3.1, Proposition~3.2 and
Lemmas~3.7--3.8, pp.~24--28.
\url{https://arxiv.org/pdf/1006.0524}.
An endpoint argument of this type already appears in the posted,
unrefereed note \emph{Boundary traces and support variation for
localized Weil operators}, by the handle \texttt{bramble\_trace},
27 October 2026, \S3.3, equations~(3.11)--(3.19).
The proof below checks the operator domains and limiting integrals
directly from the preceding classical inputs; it does not use that
note as a black-box theorem or claim novelty for its trace argument.
\url{https://rstation.ai/problems/riemann_hypothesis/papers/67}.

\begin{theorem}[AC-S37: an actual square-root logarithmic endpoint trace]
\label{AC-S37}
Let \(I=(-L,L)\), \(L>0\), and let \(A_L\) be the actual closed
Weil-window operator.  Every eigenfunction \(u\) has finite limits
\[
 \tau_-(u)=\lim_{d\downarrow0}\sqrt{\log(1/d)}\,u(-L+d),
 \qquad
 \tau_+(u)=\lim_{d\downarrow0}\sqrt{\log(1/d)}\,u(L-d).
                                                        \tag{RN.1}
\]
These are linear functionals on each eigenspace.  They may a priori
vanish; existence alone is not a normalization.

More generally, put \(G=\tfrac12\log(1+D^2)\), and let \(G_I\)
be its zero-extension part operator.  Both limits exist for
\[
 \mathcal X_I=\{u\in\operatorname{Dom}(G_I):
                      u\in L^\infty(I),\ G_Iu\in L^\infty(I)\}.
                                                        \tag{RN.2}
\]
They are bounded for the norm
\(\|u\|_\infty+\|G_Iu\|_\infty\).
For every \(t>1/2\), the maps
\(f\mapsto\tau_\pm(e^{-tA_L}f)\) are bounded linear functionals
on \(L^2(I)\).  All these assertions are unconditional.
\end{theorem}
\begin{proof}
\proofstep{1. The positive reference operator.}
The Levy density of \(G\) is
\[
 n(h)=\frac{e^{-|h|}}{2|h|},\qquad
 \int_{\mathbb R}n(h)(1-\cos(\xi h))\,dh
                      =\tfrac12\log(1+\xi^2).
                                                        \tag{RN.3}
\]
Let \(\mu\) and \(a_0\) be as in (EP.4)--(EP.6), and put
\(k=\mu-n\).  The difference is integrable: it is bounded at zero
and decays exponentially at infinity.  If \(\kappa=\int k\),
then on the whole line
\[
 H_\mu=G+\kappa I-k*,\qquad
 A_L=G_I+B_L,\qquad
 B_L=(a_0+\kappa)I-P_I(k*)P_I+R_L.
                                                        \tag{RN.4}
\]
Here \(B_L\) is bounded on both \(L^2(I)\) and \(L^\infty(I)\).
These are equalities of closed forms and of their associated
operators, with the same operator domains.  Every eigenfunction
therefore belongs to \(\mathcal X_I\), by AC-S36 and
\(G_Iu=\alpha u-B_Lu\).

Every vector in \(\mathcal X_I\) has a continuous zero extension.
Indeed \(G_I-\tfrac12L_{\Delta,I}\) is compression of the kernel
\[
                   q(h)=\frac{1-e^{-|h|}}{2|h|},\qquad q(0)=\tfrac12.
                                                        \tag{RN.5}
\]
To verify this, integrate the inverse transform of
\(b/(b^2+\xi^2)\) over \(0<b<1\).
The resulting multiplier is \(\tfrac12\log(1+\xi^{-2})\).
The compressed kernel is bounded on a finite interval, on both
\(L^2\) and \(L^\infty\).  Consequently the argument in step~4
of AC-S36 and the same Chen--Weth weak-solution regularity theorem
apply to \(u\in\mathcal X_I\).

\proofstep{2. Renewal normalization.}
For the symbol \(\Psi(\xi)=\tfrac12\log(1+\xi^2)\), normalize the
renewal function \(V\) by
\[
 \mathfrak k(s)=\exp\left\{\frac1\pi\int_0^\infty
                    \frac{s\log\Psi(r)}{s^2+r^2}\,dr\right\},
 \qquad
 \int_0^\infty e^{-sy}\,dV(y)=\mathfrak k(s)^{-1}.
                                                        \tag{RN.6}
\]
The cited primary results supply \(V(0)=0\), a positive, continuous,
decreasing density \(v=V'\), and the Green kernel of the part
operator \(G_+\) on \((0,\infty)\):
\[
 \mathcal G_+(x,y)
       =\int_0^{\min(x,y)}v(x-z)v(y-z)\,dz.
                                                        \tag{RN.7}
\]
The source uses \(\log(1+\xi^2)\).  Multiplying its time generator
by \(1/2\) divides its ladder exponent by \(\sqrt2\), multiplies
\(V,v\) by \(\sqrt2\), and multiplies the Green kernel by two.
Thus (RN.6)--(RN.7) have exactly the displayed constants.
Fubini and positivity give the useful elementary bound
\[
 \int_0^R\mathcal G_+(x,y)\,dy
  =\int_0^{\min(x,R)}v(x-z)V(R-z)\,dz
  \leq V(x)V(R).
                                                        \tag{RN.8}
\]

We verify the leading constant instead of using comparability.
Substitute \(r=st\) in (RN.6).  Since
\(\int_0^\infty(1+t^2)^{-1}dt=\pi/2\),
\[
 \log\frac{\mathfrak k(s)}{\sqrt{\Psi(s)}}
   =\frac1\pi\int_0^\infty
       \frac{\log(\Psi(st)/\Psi(s))}{1+t^2}\,dt\longrightarrow0.
                                                        \tag{RN.9}
\]
The logarithm tends to zero for every \(t>0\); its absolute value
is at most \(2|\log t|\), by monotonicity and concavity of
\(r\mapsto\log(1+r)\).  This is an integrable majorant.
Thus \(F(s):=\mathfrak k(s)^{-1}\sim(\log s)^{-1/2}\).
For \(\epsilon,M>0\), monotonicity of \(V\) gives
\[
 V(x)\leq e^\epsilon F(\epsilon/x),\qquad
 V(x)\geq F(M/x)-e^{-M/2}F(M/(2x)).
                                                        \tag{RN.10}
\]
The second inequality follows by splitting the Laplace integral
at \(x\) and bounding its tail by
\(e^{-M/2}F(M/(2x))\).
Divide by \(F(1/x)\), let \(x\downarrow0\), and then let
\(\epsilon\downarrow0\), \(M\to\infty\).  We obtain
\[
                        V(x)\sqrt{\log(1/x)}\longrightarrow1.
                                                        \tag{RN.11}
\]

\proofstep{3. The operator domain at the other endpoint.}
Translate the interval to \((0,\ell)\), where \(\ell=2L\),
and write \(h=G_{(0,\ell)}u\).  The cross operator from
\((0,\ell)\) to \((\ell,\infty)\) is
\[
 (Cu)(y)=\int_0^\ell n(y-z)u(z)\,dz.
\]
It has \(L^2\) norm at most \(\pi/2\).
Indeed, after measuring distance from \(\ell\), its kernel is
bounded by \(1/(2(s+t))\).  Schur's test with the weight
\(t^{-1/2}\) uses the exact integral
\(\int_0^\infty t^{-1/2}(s+t)^{-1}dt=\pi s^{-1/2}\).

For completeness, this bound also justifies splitting the form
domain.  Truncate \(n\) by \(1_{|h|>\epsilon}\).
The form of \(f=f_1+f_2\), with disjoint supports on the two
intervals, is the sum of the two part forms minus
\(2\operatorname{Re}\langle C_\epsilon f_1,f_2\rangle\).
The cross norms have the preceding uniform bound.
Monotone convergence of the nonnegative part forms shows that
\(f\) is in the half-line form domain if and only if its two
restrictions are in their respective part domains.
The cross operators converge strongly, first on functions supported
away from the splitting point and then by their uniform bound.
Consequently zero extension of \(u\in\operatorname{Dom}(G_{(0,\ell)})\)
belongs to \(\operatorname{Dom}(G_+)\), with
\[
 G_+u=g,\qquad
 g(y)=\begin{cases}
 h(y),&0<y<\ell,\\
 -\displaystyle\int_0^\ell n(y-z)u(z)\,dz,&y>\ell.
 \end{cases}
                                                        \tag{RN.12}
\]
This follows by testing the split form against both components.
It establishes the operator identity without an unproved boundary
term.  For bounded \(u,h\), the exterior part is
\(O(1+|\log(y-\ell)|)\) near \(\ell\) and decays exponentially
at infinity.  Hence \(g\in L^1\cap L^2(0,\infty)\).

\proofstep{4. Absolute Green representation and the trace.}
The kernel of \(G_+\) in \(L^2\) is zero: zero form energy forces
the Fourier transform of the zero extension to vanish off \(0\),
which forces the vector to be zero.  By the spectral theorem,
\[
 \int_0^T e^{-sG_+}g\,ds
          =u-e^{-TG_+}u\longrightarrow u\quad\hbox{in }L^2.
                                                        \tag{RN.13}
\]
For \(x<\ell/4\), the Green potential of \(|g|\) is finite.
Near zero this follows from (RN.8) and boundedness of \(h\).
For \(y\geq\epsilon\) and \(x<\epsilon/2\), monotonicity gives
\[
 \frac{\mathcal G_+(x,y)}{V(x)}
              \leq v(y-x)\leq v(y/2)\leq v(\epsilon/2).
                                                        \tag{RN.14}
\]
Since \(g\in L^1\), the positive semigroup kernel, (RN.13),
and absolute convergence identify \(u\) almost everywhere near
zero with \(\int_0^\infty\mathcal G_+(x,y)g(y)dy\).
For each fixed \(y>0\), the quotient in (RN.14) tends to \(v(y)\),
by continuity of \(v\) and (RN.7).
The contribution from \(0<y<\epsilon\), divided by \(V(x)\),
is at most \(\|h\|_\infty V(\epsilon)\) when \(x<\epsilon/2\).
Its proposed limit has the same bound.  Dominated convergence
on the complement, followed by \(\epsilon\downarrow0\), proves
\[
 \tau_-(u)=\int_0^\ell v(y)h(y)\,dy
 -\int_0^\ell u(z)\int_\ell^\infty v(y)n(y-z)\,dy\,dz.
                                                        \tag{RN.15}
\]
Here \(u,h\) have been translated with the interval.
Equation (RN.11) fixes the coefficient in (RN.1).
The estimate holds uniformly for sufficiently small \(x\) outside
a null set.  Continuity of \(u\) and \(V\) extends it to every
interior point; hence the limit is an ordinary one-sided limit.
Reflection gives \(\tau_+\).
Both integrals in (RN.15) are absolute, and
\[
 |\tau_\pm(u)|
    \leq V(\ell)\|G_Iu\|_\infty
                    +\tfrac12v(\ell)\|u\|_\infty.
                                                        \tag{RN.16}
\]
For the second term use
\(\int_0^\ell\int_\ell^\infty n(y-z)dy\,dz
 =\int_0^\infty\min(t,\ell)n(t)dt\leq1/2\).

\proofstep{5. A bounded functional after smoothing.}
AC-S36 gives \(e^{-rA_L}:L^2\to L^\infty\) for \(r>1/2\).
Given \(t>1/2\), choose \(1/2<r<t\).
Semiboundedness and the spectral theorem make
\(A_Le^{-(t-r)A_L}\) bounded on \(L^2\), so
\[
 A_Le^{-tA_L}=e^{-rA_L}A_Le^{-(t-r)A_L}:L^2\to L^\infty
\]
is bounded.  Equation (RN.4) then bounds
\(G_Ie^{-tA_L}:L^2\to L^\infty\).
Applying (RN.16) proves the last assertion.
\end{proof}

\begin{proposition}[AC-S38: nonzero trace on the proved positive-kernel range]
\label{AC-S38}
Let \(r_0>1\) be the root of \(r_0^3-r_0-1=0\), and suppose
\(0<L\leq\tfrac12\log r_0\).
The nonnegative, simple, even lowest eigenfunction furnished by AC-S28
has
\[
                         \tau_-(u)=\tau_+(u)>0.
                                                        \tag{RN.17}
\]
It therefore has a unique positive scalar normalization
\(\tau_+(u)=1\).  This conclusion is unconditional.
\end{proposition}
\begin{proof}
There are no prime translations in this range, since
\(e^{2L}\leq r_0<2\).  On the whole line define the nonnegative
symmetric jump density
\[
 j_L(h)=\mu(h)-2\cosh(h/2)1_{|h|\leq2L}.
                                                        \tag{RN.18}
\]
Its nonnegativity and strict positivity for \(0<|h|<2L\)
are exactly the kernel inequality proved in AC-S28.
It has the same leading singularity \(1/(2|h|)\).
If \(H_j\) is its jump operator and
\(M_p=\int_{-2L}^{2L}2\cosh(h/2)dh\), then
\[
                         A_L=(H_j)_I+(a_0+M_p)I.
                                                        \tag{RN.19}
\]
The part jump form is strictly positive: the mass of jumps longer
than \(2L\), for which the support and its translate are disjoint,
already supplies a positive multiple of \(\|u\|_2^2\).
Thus the nonnegative ground state satisfies
\((H_j)_Iu=\beta u\geq0\), with \(\beta>0\).
Translate \(I\) to \((0,\ell)\).

We use a compact nonnegative barrier \(b\) which equals
\((\log(2/x))^{-1/2}\) near \(0+\), is zero for \(x\leq0\),
and is smoothly cut off before \(2\).
This is the one-dimensional barrier of
Hern\'andez-Santamar\'ia, L\'opez R\'ios and Salda\~na,
arXiv:2401.18033v2, (2.5)--(2.6), Theorem~2.4,
and its complete proof in \S2.1, pp.~8--11.
\url{https://arxiv.org/pdf/2401.18033}.
Here are the estimates needed for this application.
The barrier is bounded with bounded variation, hence belongs to
\(H^s(\mathbb R)\) for every \(0<s<1/2\), by
\(\|b(\cdot+h)-b\|_2^2\leq2\|b\|_\infty
\operatorname{Var}(b)|h|\).
In particular it belongs to the logarithmic form domain.
Its logarithmic Laplacian is bounded in \((0,2)\).
Away from zero this follows from smoothness.  At \(x=\epsilon\),
split the singular integral over
\((\epsilon-1,0)\), \((0,2\epsilon)\), and
\((2\epsilon,1+\epsilon)\), taking the cutoff equal to one
on a neighborhood of \([0,1]\).
The middle difference integral is \(o(1)\), by splitting once
more at \(\epsilon/2\) and using
\(b'(x)=O(x^{-1}\log(2/x)^{-3/2})\).
The other two terms sum to
\[
 2b(\epsilon)\log(1/\epsilon)
 -\int_\epsilon^1\frac{dt}{t\sqrt{\log(2/(t+\epsilon))}}
                     =2\sqrt{\log2}+o(1).
                                                        \tag{RN.20}
\]
To justify replacing \(t+\epsilon\) by \(t\), put
\(z=-\log t\).  The difference tends pointwise to zero and,
on \(1\leq z\leq\log(1/\epsilon)\), is at most
\(Cz^{-3/2}\), because
\(e^{-z}+\epsilon\leq2e^{-z}\) and
\(\epsilon e^z\leq1\).
It is uniformly bounded on \(0<z<1\), and is extended by zero
beyond \(\log(1/\epsilon)\).
Dominated convergence proves the claimed \(o(1)\).
The remaining nonsingular integral and scalar term are bounded.
This verifies the barrier input with an integrable majorant.

Continuity, nonnegativity and nontriviality of \(u\) give a compact
interval \(B\subset(0,\ell)\) on which \(u\geq c_B>0\).
Choose \(\delta>0\) with \(3\delta<\inf B\), put
\(b_\delta(x)=b(x/\delta)\), and let \(D=(0,2\delta)\).
The barrier is supported strictly before the right endpoint of \(D\).
Scaling gives
\(L_\Delta b_\delta(x)=(L_\Delta b)(x/\delta)
 -2\log\delta\,b_\delta(x)\).
The difference between \(H_j\) and \(\tfrac12L_\Delta\),
when applied to compactly supported functions on this fixed interval,
is bounded on \(L^\infty\), by the same kernel subtraction
as (EP.14).  Hence \(H_jb_\delta\) is bounded on \(D\).
Strict positivity and continuity of \(j_L\) at the separated
interior distances give
\[
 \inf_{x\in D}\int_Bj_L(x-y)dy=:\kappa_D>0,
 \qquad H_j1_B(x)=-\int_Bj_L(x-y)dy\leq-\kappa_D.
                                                        \tag{RN.21}
\]
Choose \(C>0\) so that
\(\phi=b_\delta+C1_B\) satisfies \(H_j\phi\leq0\) on \(D\).
Outside \(D\), the function \(\phi\) equals \(C1_B\).
The choice of \(B\) therefore gives
\(u-\varepsilon\phi\geq0\) off \(D\) for some \(\varepsilon>0\).
It is also a weak supersolution of \(H_j\) in \(D\).

The required comparison follows directly from the nonnegative jump
form: test against the negative part of
\(u-\varepsilon\phi\), which is supported in \(D\).
The pairing is at most minus the energy of that negative part,
and is at the same time nonnegative.
Strict positivity of the part form forces the negative part to vanish.
All test functions belong to the closed jump-form domain;
truncation and form-core approximation justify this test.
We have proved \(u\geq\varepsilon b_\delta\) near zero.
Since \(\sqrt{\log(1/x)}b_\delta(x)\to1\), AC-S37 gives
\(\tau_-(u)\geq\varepsilon>0\).
Reflection and evenness give (RN.17).
\end{proof}

\begin{corollary}[AC-S39: convergent Galerkin normalization by the actual trace]
\label{AC-S39}
Fix a window for which the lowest eigenvalue \(\alpha\) is simple
and even, and its normalized eigenvector \(u\) has \(\tau_+(u)\ne0\).
These hypotheses hold throughout the range of AC-S38.
Let \(u_N\) and \(\alpha_N\) be the normalized lowest Galerkin
vectors and eigenvalues, with phases chosen as in AC-S06 and its
semibounded extension following AC-S28.
For any fixed \(t>1/2\), set
\[
 b_N=e^{t\alpha_N}\tau_+(e^{-tA_L}u_N).
                                                        \tag{RN.22}
\]
Then \(b_N\to\tau_+(u)\ne0\), and
\[
 \frac{u_N}{b_N}\longrightarrow\frac{u}{\tau_+(u)}
 \quad\hbox{in }L^2(I)\hbox{ and in form norm}.
                                                        \tag{RN.23}
\]
Their Fourier transforms converge locally uniformly in the whole
complex plane.  The finite regularized determinants in S25,
Theorem~5.10 (\texttt{finmain}), admit the corresponding nonzero
scalar renormalization with this limit.
\end{corollary}
\begin{proof}
AC-S37 makes \(f\mapsto\tau_+(e^{-tA_L}f)\) a bounded
\(L^2\) functional.  Since \(u_N\to u\) and
\(\alpha_N\to\alpha\),
\[
 b_N\longrightarrow e^{t\alpha}\tau_+(e^{-tA_L}u)
                  =\tau_+(u).
\]
Division by these nonzero limiting scalars preserves both the
\(L^2\) and form-norm convergence of AC-S06.
Compact support gives the locally uniform entire Fourier convergence
by the same Cauchy--Schwarz estimate as in that proposition.
The source determinant identity is valid for the eventually
simple even finite matrices; multiplying its source normalization
by the ratio of its finite boundary value to \(b_N\) gives exactly
the transform of \(u_N/b_N\), with the fixed source phase removed
as in AC-S06.  This is an explicit scalar renormalization of that
identity, not an assumption about a limiting point evaluation.
\end{proof}

\paragraph{Precise scope of the repair.}
Equation (RN.1) defines a genuine continuum endpoint functional,
and AC-S38 proves that it normalizes the ground state on a specified
nonempty range.  Equation (RN.22) gives a continuous way of applying
that functional to finite Galerkin vectors; it uses the actual
window semigroup.  A finite trigonometric vector with nonzero endpoint
value generally has no finite trace (RN.1) before this smoothing.
AC-S39 itself does not estimate that unsmoothed value.
AC-S40 below proves its decay, and AC-S41 determines the exact
nonzero logarithmic coefficient.  AC-S42 handles every fixed
simple-even window, including its prime translations.
The unrescaled source determinants diverge wherever the limiting
Fourier transform is nonzero; division by \(\sqrt{\log N}\)
gives a locally uniform entire limit.
AC-S43 identifies its coefficient exactly with the renewal trace
and proves that trace nonzero for every fixed simple-even minimum.
The simple-even lowest-line assertion on arbitrary windows and
selection of the large-window ground line remain open.
None of these fixed-window results identifies the prolate
spectrum with zeta zeros or proves the S33 gap hypothesis.
